\scanpage{50}
% archive sheet number: 49

J'aimerais comprendre la relation entre les $\sigma_a^F$, et les signes $\varepsilon_F(a)$.

Mais il faut encore introduire une autre structure

c) L'ensemble $C_K\subset S_K$ des ``pts critiques'' avec la condition que \inserted{ils sont d'ordre 2, et que} sur chaque \deleted{\unsure{ar}} grande arête, il y a \deleted{\unsure{au plus}} \inserted{exactement} un \deleted{tel} élément de $C_{\unsure{K|F}}$ \deleted{(quand il n'y en a pas, l'arête est dite \uline{exceptionnelle}).} De plus, chaque p\unsure{ts|t} critique $s$ est muni d'un poids $\nu(s)\geq\deleted{\unsure{2}}1$, \inserted{qui est pair ou impair, suivant que $s$ est noir, ou \deleted{orange}} (le poids d'un élément de $S_K$ est par définition $\nu(s)$ si \unsure{critique}, 1 si $s$ non critique \deleted{\unsure{du pt}} et de $S'_K$, 0 s'il est non critique \deleted{et} est de $S''_K$)

\margin{définir les ``arcs critiques'' ou ``arcs de transition'' sur les arêtes, d'amplitude totale $2\nu(s)$ (il contient $\nu(s)$ sommets \deleted{de $S'_K$} rouges, et ses extrémités sont \unsure{rouges|oranges})}

% slanted
\margin{\uline{NB} Les arcs de transition (fermés) sont mutuellement disjoints}

% slanted
\margin{Les pts critiques sont les centres des arcs critiques, et déterminent les paliers (comme la moitié de l'amplitude correspondante du cas entier) --- donc il \unsure{revient précisément} au même de se donner ces arcs critiques.}

% slanted
\margin{les sommets d'ordre 4 sont exactement \deleted{\unsure{ceux qui}} les milieux des arcs \deleted{\unsure{complémentaires}} \inserted{compris} \deleted{\unsure{par}} entre les arcs \deleted{de transition} \inserted{critiques}}

En termes des sommets critiques, on a ceci :

1) Soit $s$ un sommet critique, incident à la face $F$. Soit $s'$ son antipodique dans $\partial F$. \deleted{Alors si $s'$ est non critique, on doit le signe $\varepsilon_a$} \inserted{Pour qu'on ait $\varepsilon_F(a)=+1$, il f et s.} que le sommet $s'$ antipodique \deleted{\unsure{$t$}} de $s$ \inserted{pour $F$} soit égalmt critique (il a alors le même indice)

% slanted
\margin{relation entre $C_K$ et les signes $\varepsilon_F(a)$}

2) Sur toute face $F$, il y a au moins 3 pts critiques sur $\partial F$ (et même au moins 4, \deleted{\unsure{3)} si tous les pts} il y a au moins un pt critique \inserted{$s$} de signe $+1$. \deleted{\unsure{aura faites}} \inserted{Dans ce cas,} sur chacune
